% Md Shibly Sadique - one-page resume
% Build: pdflatex resume.tex   (run twice for stable links)
% Companion to cv.tex, which is the long-form academic version.
% Typeface and metrics follow the reference template: Libertinus Serif, 10pt,
% 0.6in side margins, centered header, italic summary between rules.
\documentclass[10pt,letterpaper]{article}

\usepackage[left=0.58in,right=0.58in,top=0.4in,bottom=0.35in]{geometry}
\usepackage{libertinus}
\usepackage{microtype}
\usepackage{enumitem}
\usepackage{xcolor}
\usepackage[hidelinks]{hyperref}

\definecolor{link}{HTML}{1F4E8C}
\definecolor{muted}{HTML}{3A3A3A}

\hypersetup{
  colorlinks=true, urlcolor=link, linkcolor=link,
  pdfauthor={Md Shibly Sadique},
  pdftitle={Resume - Md Shibly Sadique}
}

\setlength{\parindent}{0pt}
\setlength{\parskip}{0pt}
\linespread{0.99}
\pagestyle{empty}
\raggedright

\SetTracking{encoding=*}{110}
\newcommand{\spaced}[1]{\lsstyle #1}

% section heading: letterspaced small caps with a rule running to the right margin
\newcommand{\rsection}[1]{%
  \vspace{4.5pt}%
  \noindent{\scshape\spaced{#1}}\hspace{6pt}%
  {\leaders\hrule height \dimexpr0.4ex+0.9pt\relax depth -0.4ex\hfill}\kern0pt%
  \par\vspace{6pt}}

\setlist[itemize]{leftmargin=12pt,itemsep=0pt,topsep=1.5pt,parsep=0pt,label=\textbullet}

% \job{employer}{role}{location}{dates}
\newcommand{\job}[4]{%
  \noindent{\scshape #1} \,-\, \textit{#2}\hfill{\textit{#3}\quad\textbf{#4}}\par\vspace{1pt}}

\begin{document}

%%% ---------------------------------------------- header
\begin{center}
{\LARGE\scshape\spaced{Md Shibly Sadique}}\par\vspace{2pt}
PhD Candidate, Biomedical Image Analysis, Old Dominion University\par\vspace{2pt}
{\small +1 (757) 895-6050 \;$\cdot$\;
\href{mailto:msadi002@odu.edu}{msadi002@odu.edu} \;$\cdot$\;
\href{https://shiblyg.github.io}{shiblyg.github.io} \;$\cdot$\;
\href{https://www.linkedin.com/in/md-shibly-sadique-277b891a5/}{linkedin.com/in/md-shibly-sadique} \;$\cdot$\;
\href{https://scholar.google.com/citations?user=EWnmhbgAAAAJ&hl=en}{Google Scholar}}
\end{center}

\vspace{7pt}

\textit{PhD candidate with eight years of research experience in medical image analysis and machine learning, graduating July
2027. Demonstrated expertise in designing, training, and evaluating deep learning models with PyTorch, TensorFlow, and
Scikit-learn for 3D volumetric MRI and CT. Authored fifteen peer-reviewed publications and led a team in an international
segmentation challenge. Aiming to apply uncertainty-aware modeling skills to deliver reliable clinical AI systems.}

\vspace{3pt}

%%% ---------------------------------------------- experience
\rsection{Experience}

\job{Old Dominion University}{Graduate Research Assistant, Vision Lab}{Norfolk, VA}{2018 - Present}
\begin{itemize}
\item Developed UNCE-NODE, a neural ODE-enhanced U-Net for multimodal MRI tumor segmentation, cutting trainable parameters by roughly 30\% and accelerating inference while holding Dice accuracy.
\item Built a classifier separating recurrent brain tumor from radiation necrosis on multiresolution radiomic features, reaching AUC $0.892 \pm 0.055$ on NIH/NIBIB-funded data.
\item Designed a copula-based survival pipeline handling dependent censoring in glioblastoma; cross-validated AUC $\approx 0.77$, outperforming Cox regression.
\item Implemented an uncertainty quantification framework over ensemble nnU-Nets, combining entropy and mutual-information estimates to stratify erroneous segmentations; led a team to eighth place in the AMOS 2022 challenge.
\end{itemize}

\vspace{2pt}
\job{Old Dominion University}{Graduate Teaching Assistant}{Norfolk, VA}{2019 - Present}
\begin{itemize}
\item Ran problem sessions and laboratories for circuits, linear systems, and signals courses; built MATLAB/Simulink and Multisim exercises tying theory to application.
\item Supported ECE 481W/487 Senior Design from 2021, developing course materials and ABET outcome reports.
\item Mentored two NSF-REU undergraduates through summer projects on privacy-preserving and federated medical image analysis.
\end{itemize}

\vspace{2pt}
\job{University of Dhaka}{Research Assistant, Microwave \& Optical Fiber Lab}{Dhaka, BD}{2014 - 2016}
\begin{itemize}
\item Built the lab's simulation, fabrication, and measurement capability for microstrip patch antenna research.
\end{itemize}

%%% ---------------------------------------------- publications
\rsection{Selected Publications}
\begin{itemize}
\item \textbf{Sadique, M.S.}, et al. Brain Tumor Recurrence vs. Radiation Necrosis Classification and Patient Survivability Prediction. \textit{IEEE Journal of Biomedical and Health Informatics}, 2024.
\item Temtam, A., \dots \textbf{Sadique, M.S.}, et al. Functional Brain Network Identification in Opioid Use Disorder Using Machine Learning of Resting-State fMRI. \textit{Computers in Biology and Medicine}, 2025.
\item \textbf{Sadique, M.S.}, et al. Class Activation Mapping and Uncertainty Estimation in Multi-Organ Segmentation. \textit{SPIE Medical Imaging}, 2023.
\item Twelve further peer-reviewed papers across MICCAI Brainlesion, SPIE Medical Imaging, and Springer book chapters.
\end{itemize}

%%% ---------------------------------------------- projects
\rsection{Selected Projects}
\begin{itemize}
\item \textbf{Adult-to-Pediatric Domain Adaptation} (2026) \,-\, gradient-reversal domain alignment with multiresolution fractal features, transferring 1{,}251 adult glioma cases to 228 pediatric cases; presented at SPIE Optics + Photonics.
\item \textbf{Learnability Under Sampling Noise} (2025-) \,-\, Resolvable Expressive Capacity framework with eigentask decomposition, quantifying which tumor growth directions a PDE model recovers from noisy imaging.
\end{itemize}

%%% ---------------------------------------------- skills
\rsection{Technical Skills}
\begin{itemize}
\item \textbf{Languages \& tools:} Python, MATLAB, R, C, Shell, Git, \LaTeX, Docker, CUDA, Slurm (HPC)
\item \textbf{Machine learning:} PyTorch, TensorFlow, Keras, Scikit-learn, OpenCV; CNN, RNN, LSTM, U-Net, nnU-Net, Neural ODEs, Transformers, Diffusion Models, Transfer Learning, Domain Adaptation
\item \textbf{Medical imaging:} FSL, FreeSurfer, 3D Slicer, ANTs, SimpleITK, NiBabel, ITK-SNAP, MIPAV, MedPerf
\item \textbf{Statistics:} Survival analysis, copula modeling, uncertainty quantification, feature selection, PCA, LDA
\end{itemize}

%%% ---------------------------------------------- education
\rsection{Education}

\noindent\textbf{Ph.D., Electrical \& Computer Engineering} \,-\, Old Dominion University\hfill{\small 2018 - July 2027 (expected)}\par
\noindent\textbf{M.Sc., Electrical \& Electronic Engineering} \,-\, University of Dhaka\hfill{\small 2014 - 2016}\par
\noindent\textbf{B.Sc., Electrical \& Electronic Engineering} \,-\, University of Dhaka\hfill{\small 2010 - 2014}\par

%%% ---------------------------------------------- honors
\rsection{Honors}
SPIE Medical Imaging Student Conference Support (2022, 2023) \;$\cdot$\; Team Lead, Eighth Place, AMOS Multi-Organ
Segmentation Challenge (2022) \;$\cdot$\; BraTS Challenge participant (2021-2026) \;$\cdot$\; NSICT Fellowship for M.S.
Thesis, Bangladesh (2015) \;$\cdot$\; IEEE Graduate Student Member; SPIE Student Member

\end{document}
